\chapter{Task YAML Structure}
\label{ch:task-yaml}

\newthought{A task card is one \textsc{yaml} file} that describes a whole scan. This chapter is
the map: it shows the handful of top-level blocks a card can contain, says which are required,
and points you to the chapter that documents each one in detail. Read it once to get the shape;
then dip into the later chapters as you write.

\section{The shape of a card}
\label{sec:card-shape}

Recall the three-step model from Chapter~\ref{ch:introduction}: for every point, \Jarvis\
\emph{draws a point}, \emph{runs a workflow} to get observables, and \emph{scores} them with a
likelihood. A task card simply fills in those steps. Here is the full skeleton---you will rarely
use every block at once:

\begin{yamlwin}
Scan:          # run name + where outputs go          (required)
Sampling:      # how points are drawn + the objective  (required)
EnvReqs:       # platform + dependency checks           (recommended)
Calculators:   # external-program workflow  )  choose
Operas:        # in-process workflow        )  ONE of these
LibDeps:       # shared external libraries             (optional)
Utils:         # helper functions for expressions      (optional)
\end{yamlwin}

\section{Required, recommended, optional}
\label{sec:required-optional}

Only two blocks are mandatory; the rest you add when you need them
(Table~\ref{tab:blocks}).

\begin{table}[ht]
  \footnotesize
  \begin{tabular}{@{}llp{0.46\linewidth}@{}}
    \toprule
    \textbf{Block} & \textbf{Status} & \textbf{What it does} \\
    \midrule
    \yamlkey{Scan}        & required    & names the run and sets the output location \\
    \yamlkey{Sampling}    & required    & chooses the sampler, the scan variables, and the likelihood \\
    \yamlkey{EnvReqs}     & recommended & checks the machine before the run starts \\
    \yamlkey{Calculators} & one backend & runs external programs to produce observables \\
    \yamlkey{Operas}      & one backend & runs in-process Python operators instead \\
    \yamlkey{LibDeps}     & optional    & installs shared external libraries once \\
    \yamlkey{Utils}       & optional    & defines helper functions used in expressions \\
    \bottomrule
  \end{tabular}
  \caption{The top-level blocks of a task card.}
  \label{tab:blocks}
\end{table}

\begin{jwarn}
Pick \emph{one} workflow backend: a card uses \yamlkey{Calculators} \emph{or} \yamlkey{Operas},
not both. Use \yamlkey{Operas} for quick, in-process calculations and \yamlkey{Calculators} to
drive external programs (Chapter~\ref{ch:core-concepts} explains the choice).
\end{jwarn}

\section{A complete minimal card}
\label{sec:minimal-card}

The quick-start from Chapter~\ref{ch:quickstart} is a real, complete card. Stripped to its
essentials it is just \yamlkey{Scan} + \yamlkey{Sampling} + a backend:

\begin{yamlwin}
Scan:
  name: "my_first_scan"
  save_dir: "&J/outputs"

Sampling:
  Method: "Random"
  Variables:
    - name: x
      distribution: {type: Flat, parameters: {min: 0, max: 5}}
  LogLikelihood:
    - {name: "LogL", expression: "-0.5 * (x - 2)**2"}

Operas:
  make_paraller: 4
  Modules: []
\end{yamlwin}

\noindent Everything else in this part of the book is about filling these blocks in with more
detail and confidence.

\section{Where each block is documented}
\label{sec:block-map}

\begin{description}[style=nextline, leftmargin=1.2em]
  \item[\yamlkey{Scan}] Chapter~\ref{ch:scan-block} --- run name, output directory, sample
        retention.
  \item[\yamlkey{Sampling}] Chapter~\ref{ch:sampling-block} (the block),
        Chapter~\ref{ch:variables} (scan variables), Chapter~\ref{ch:loglikelihood}
        (the likelihood), and Part~\ref{part:samplers} (one chapter per sampler).
  \item[\yamlkey{EnvReqs}] Chapter~\ref{ch:envreqs} --- platform and dependency checks.
  \item[\yamlkey{Calculators}] Chapter~\ref{ch:calculators} --- external-program workflows.
  \item[\yamlkey{Operas}] Chapter~\ref{ch:operas} --- in-process operator workflows.
  \item[\yamlkey{LibDeps}] Chapter~\ref{ch:libdeps} --- shared external libraries.
  \item[\yamlkey{Utils}] Chapter~\ref{ch:utils} --- helper functions for expressions.
\end{description}

\noindent Two cross-cutting topics support all of these: the \emph{expressions} you write in
likelihoods and mappings (Chapter~\ref{ch:symbols}) and the \emph{path markers and tokens}
that keep paths portable (Chapter~\ref{ch:core-concepts}).

\begin{jtip}
\textsc{yaml} is whitespace-sensitive: indent with spaces (never tabs), and keep items under a
key aligned. If \Jarvis\ reports a validation error, it almost always points at the block and
key where the indentation or a value is wrong.
\end{jtip}
